%% ch05 --- Funkcje i sterowanie: dekoratory, składanie list, generatory, menedżery kontekstu
%%
%% Napisane we wrześniu 2026. Każdy listing w code/ch05/ uruchamia się na
%% przypiętym interpreterze i jest wykonywany przez code/tests/test_listings.py;
%% każde cytowane wyjście to wygenerowany transkrypt; jedyny zmierzony wynik
%% pochodzi z code/measure/ch05_concat.py i trafia na stronę jako ograniczenie.

\chapter{Funkcje i sterowanie: dekoratory, składanie list, generatory, menedżery kontekstu}
\label{ch:functions}

\noindent
Słownik sterowania w Pythonie odwzorowuje się na C\# funkcja po funkcji.
Dekorator to middleware, składanie listy to zapytanie LINQ, generator to
blok iteratora, \code{with} to \code{using}, \code{match} to wyrażenie
switch. To odwzorowanie jest dość bliskie, by warto je było zrobić, i dość
luźne, by warto je było sprawdzić. Ten rozdział jest o miejscach, w których
się rozjeżdża: każde z nich daje Pythona, który się wykonuje, a trzy z nich
dają Pythona, który wykonuje się po cichu źle.

Argument mieści się w jednym zdaniu. \textbf{Dekorator się wykonuje;
atrybut nie}, a prawie wszystko inne w tym rozdziale to ta sama różnica w
innym przebraniu \dash{} co już się wydarzyło, zanim dostałeś wartość do
ręki, a co dopiero czeka.

\section{Funkcja jest obiektem}
\label{sec:ch05-first-class}

Znasz już ten pomysł, pod dwiema nazwami. Delegat jest typem;
\code{Func<>} jest rodziną typów indeksowaną liczbą argumentów; grupa metod
konwertuje się na jeden z nich. Python ma ten pomysł bez ceremonii: funkcja
jest zwykłą wartością zwykłego typu, można ją przechować, przekazać, zwrócić
i włożyć do słownika, a adnotacja jest jedynym miejscem, w którym sygnatura
jest zapisana.

\pyregion{code/ch05/first_class.py}{table}{Tablica wysyłkowa: bez deklaracji delegata i bez liczby argumentów w nazwie typu}{lst:ch05-table}

\begin{csbox}
Powyższy słownik to \code{Dictionary<string, Func<Job, bool>>}. Dwie różnice
warte zapamiętania. \code{Callable} bierze swoje parametry jako listę, więc
liczba argumentów jest kształtem, a nie częścią nazwy, i nie ma osobnego
\code{Action}: funkcja nic nie zwracająca to funkcja zwracająca \code{None}.
Nie ma też \code{Invoke} ani konwersji grupy metod, bo nigdy nie było typu
delegata, na który miałaby konwertować.
\end{csbox}

\noindent
To cała maszyneria. Dekorator to funkcja nałożona na funkcję,
\code{functools.partial} to funkcja zwracająca funkcję, a funkcja
kluczująca to funkcja przekazana innej. Nic poniżej nie wprowadza nowego
rodzaju rzeczy.

\section{Dekorator się wykonuje; atrybut nie}
\label{sec:ch05-decorators}

To nagłówek rozdziału i najdroższy nawyk, jaki inżynier .NET przynosi do
Pythona. \code{[Timed]} w C\# to metadane. Kompilator je zapisuje, nic nie
zmieniają i robią cokolwiek dopiero wtedy, gdy framework odczyta je przez
refleksję i postanowi zadziałać. Czytanie \code{@timed} w ten sam sposób to
czytanie go jako czegoś biernego, a bierny nie jest: \code{@timed} to
wywołanie funkcji, które dzieje się w chwili, gdy kończy się instrukcja
\code{def}.

\pyregion{code/ch05/decorators.py}{definition_time}{Dekorator wykonuje się w czasie definicji; ciało funkcji w czasie wywołania}{lst:ch05-deftime}

\noindent
Uruchom plik, a moduł ma już zapisaną linię, zanim cokolwiek zostanie
wywołane:

\transcript{ch05-decorators}

\mermaidfig{ch05-decorator-timeline}{Dekorator jest nakładany w czasie definicji, a to, co zwróci, jest odtąd znaczeniem nazwy}{kiedy wykonuje się dekorator}

\begin{trapbox}
\textbf{Dekorator to nie metadane.} Dostaje obiekt funkcji, a to, co zwróci,
zastępuje tę funkcję \dash{} więc dekorator może ją opakować, zalogować,
zapamiętać w cache, gdzieś zarejestrować albo zwrócić coś, co funkcją w
ogóle nie jest. Najbliższym odpowiednikiem w C\# nie jest atrybut, tylko
kawałek middleware owinięty wokół wywołania, a \code{functools.wraps}
powstrzymuje tożsamość opakowania przed wyciekaniem do każdej linii logu i
każdego stosu wywołań w dole systemu. Bez niego
\code{parse.\_\_name\_\_} to \code{wrapper}.
\end{trapbox}

\noindent
Dekorator przyjmujący argumenty to funkcja zwracająca dekorator, czyli trzy
poziomy wywołań, które czyta się gorzej, niż działają.

\pyregion{code/ch05/decorators.py}{retry}{Trzy poziomy: wywołanie, zwrócony przez nie dekorator i opakowanie}{lst:ch05-retry}

\noindent
Pomiar wywołania to druga połowa tego, co robi middleware, a złożenie obu
pokazuje kolejność, w jakiej nakładają się dekoratory. Nakładają się od
dołu, więc \code{@timed} napisane nad \code{@retry} opakowuje ponawiające
opakowanie i dzięki temu mierzy wszystkie próby, a nie jedną z nich.

\pyregion{code/ch05/decorators.py}{timing}{Druga połowa middleware i powód, dla którego nie drukuje czasu}{lst:ch05-timed}


Biblioteka standardowa dostarcza kilka, których raczej użyjesz, niż je
napiszesz. \api{functools.lru\_cache} i \api{functools.cache} zapamiętują
wynik po krotce argumentów, co znaczy, że argumenty muszą być hashowalne, a
argument typu \code{list} podnosi \code{TypeError}; obie doczepiają do
funkcji \code{cache\_info()}, czyli dokładnie to, co potrafi zrobić tylko
dekorator, który funkcję \emph{zastępuje}.

\begin{note}
Middleware agentowe jest najbardziej widocznym zastosowaniem tego mechanizmu
w pythonowym stosie AI i należy do tomu towarzyszącego: \emph{LangChain,
LangGraph and Async Python} poświęca rozdział kolejności middleware i
haczykom, które je składają. Ten rozdział ma mechanizm, tamten ma
zastosowanie.
\end{note}

\begin{exercise}{e05_01_audit}{Dekorator, który jest middleware}
Napisz \code{audit} tak, by zapisywał każde wywołanie i poza tym zostawiał
funkcję niezmienioną. Rozstrzygają cztery asercje: zwracana wartość,
zapisane wywołania oraz \code{\_\_name\_\_} i \code{\_\_doc\_\_}, które mają
przeżyć opakowanie. Te dwie ostatnie zawodzą, gdy zapomnisz o
\api{functools.wraps}, i to one są dwiema, których nikt nie zauważa, dopóki
stos wywołań trzy usługi dalej nie powie \code{wrapper}.
\end{exercise}

\section{Pięć rodzajów parametru, a C\# ma trzy}
\label{sec:ch05-parameters}

\code{*args} to \code{params}. Wartości domyślne to argumenty opcjonalne.
Przekazanie po nazwie to argument nazwany. Do tego Python ma dwa znaczniki
bez żadnego zapisu w C\# i jedną kolekcję bez odpowiednika.

\pyregion{code/ch05/parameters.py}{signature}{Wszystko, co Python potrafi powiedzieć o parametrze, w jednej sygnaturze}{lst:ch05-signature}

\noindent
Wszystko przed \code{/} można przekazać wyłącznie pozycyjnie, wszystko po
\code{*} wyłącznie po nazwie, a \code{**options} zbiera nazwy, których
funkcja nigdy nie zadeklarowała. Ten \code{/} jest wart więcej, niż wygląda.
W C\# nazwa każdego parametru jest publiczną powierzchnią na zawsze, bo
każdy wywołujący może użyć jej jako argumentu nazwanego, a zmiana nazwy go
psuje; ten znacznik jest sposobem, w jaki Python to sobie odbiera.

\pyregion{code/ch05/parameters.py}{refusals}{Dwie odmowy, których C\# nie umie wyrazić}{lst:ch05-refusals}

\transcript{ch05-parameters}

\begin{note}
Oba błędy odrzuca \pkg{pyright}, zanim cokolwiek się wykona, i dlatego
listing niesie dwa wyciszenia, żeby w ogóle dało się pokazać połowę
wykonawczą. Ostatnia linia transkryptu jest powodem, dla którego \code{/}
zarabia na swój znak: wobec parametru, który \emph{nie} jest wyłącznie
pozycyjny, ta sama pomyłka zostaje po cichu wchłonięta przez
\code{**options}, zamiast zostać odrzucona.
\end{note}

\api{functools.partial} wiąże część argumentów teraz i zwraca coś
wywoływalnego, czyli tyle, ile biblioteka standardowa ma do zaoferowania
zamiast przechwyconego delegata.

\pyregion{code/ch05/parameters.py}{partial}{partial i dodatek z 3.14, którego wtyczki typów jeszcze nie opisują}{lst:ch05-partial}

\noindent
Zmienna wartość domyślna \dash{} \code{def f(items=[])} \dash{} to druga
pułapka parametrów i należy do rozdziału~\ref{ch:typing}, bo tam właśnie
rozumowanie sprawdzarki o argumentach domyślnych gryzie pierwszy raz.

\section{Domknięcie przechwytuje zmienną, nie jej wartość}
\label{sec:ch05-closures}

Jeśli pisałeś w C\# przed wersją~5, spotkałeś ten błąd i pamiętasz lekarstwo.
Ciało \code{foreach}, które przechwytywało zmienną pętli, przechwytywało
\emph{zmienną}, więc każdy delegat widział ostatnią wartość, a kuracją było
skopiowanie jej do świeżej zmiennej lokalnej wewnątrz ciała. C\#~5 sprawił,
że \code{foreach} robi to za ciebie. Python nigdy tego nie zrobił, a
składanie listy jest pętlą.

\pyregion{code/ch05/late_binding.py}{trap}{Trzy domknięcia nad jedną zmienną}{lst:ch05-trap}

\transcript{ch05-late-binding}

\begin{trapbox}
\textbf{Każde domknięcie widzi ostatnią wartość.} Lista trzyma trzy różne
obiekty funkcji, a one dzielą jedną komórkę, więc odpowiedzią jest
\code{[2, 2, 2]}, a nie \code{[0, 1, 2]}. Obie poprawki są tą samą poprawką:
zwiąż wartość w chwili, w której funkcja powstaje \dash{} albo jako argument
domyślny, obliczany raz, przy \code{def}, albo przekazując ją do
\api{functools.partial}.
\end{trapbox}

\pyregion{code/ch05/late_binding.py}{fixes}{Świeża zmienna lokalna, wypisana wprost, i to samo przez partial}{lst:ch05-fixes}

\noindent
Tę jedną znajduje linter. \code{B023} z \pkg{ruff}, \emph{definicja funkcji
nie wiąże zmiennej pętli}, odrzuca zepsutą linię, więc powyższy listing musi
wyciszyć tę regułę, żeby w ogóle wydrukować błąd. To rzadka dobra wiadomość
w tym rozdziale: spośród tutejszych pułapek tylko tę łapie za ciebie zestaw
narzędzi z rozdziału~\ref{ch:packaging}.

\begin{exercise}{e05_02_handlers}{Zwiąż zmienną pętli}
\code{broken} zostaje zepsute \dash{} test to sprawdza, bo cała rzecz polega
na tym, że to się kompiluje i wykonuje. Napisz obie poprawki. Czwarty test
woła jeden z handlerów dwa razy i sprawdza, czy odpowiada tak samo, i to jest
ta asercja, która pada, gdy sięgniesz po współdzieloną strukturę mutowalną.
\end{exercise}

\section{Składanie listy to nie zapytanie LINQ}
\label{sec:ch05-comprehensions}

Składanie listy czyta się jak \code{Select} i \code{Where} z przestawionymi
klauzulami, a przestawienie nie jest tą różnicą, która ma znaczenie.

\begin{csbox}
\begin{csharp}
var worst = jobs
    .Where(j => j.Status == "failed")
    .OrderByDescending(j => j.Seconds)
    .Select(j => $"{j.Name} ({j.Seconds}s)")
    .Take(2);
\end{csharp}
Nic się nie wykonało. \code{worst} jest typu \code{IEnumerable<string>}, a
przejście po nim dwa razy uruchamia cały potok dwa razy.
\end{csbox}

\pyregion{code/ch05/pipeline.py}{eager}{Ten sam potok jako składanie listy. Gdy wraca, już się wykonał}{lst:ch05-eager}

\noindent
Zamień nawiasy kwadratowe na okrągłe, a stanie się odroczony \dash{} i
wtedy różni się od \code{IEnumerable} w sposób, który gryzie.

\pyregion{code/ch05/pipeline.py}{deferred}{Odroczony i zużyty przez pierwsze przejście}{lst:ch05-deferred}

\transcript{ch05-pipeline}

\mermaidfig{ch05-eager-vs-lazy}{Trzy zapisy jednego potoku i to, co każdy z nich zdążył zrobić, zanim wrócił}{zachłanny, odroczony i powtarzalny}

\begin{trapbox}
\textbf{Wyrażenie generatora jest jednoprzebiegowe.} Zapytanie LINQ jest
odroczone i powtarzalne: przejdź po nim jeszcze raz, a potok pobiegnie
jeszcze raz. Wyrażenie generatora jest odroczone i \emph{zużywalne} \dash{}
drugie przejście jest puste, bez błędu i bez ostrzeżenia. Odroczony potok
nie jest więc zamiennikiem \code{IEnumerable}, które miałeś komuś oddać, a
lekarstwem jest zmaterializować go raz, świadomie, przez \code{list()}.
\end{trapbox}

\noindent
Zauważ też, co sortowanie zrobiło z wersją leniwą: nie ma go tam.
Porządkowanie nie odpowie, dopóki nie zobaczy każdego elementu, więc potok,
który sortuje, już się wykonał, cokolwiek obiecywała reszta. Dotyczy to tak
samo \code{OrderBy}; to jedno z niewielu miejsc, w których oba języki zgadzają
się dokładnie. Gdy składanie przestaje się dobrze czytać, następną rzeczą do
sięgnięcia jest \pkg{itertools}: \api{islice} to \code{Take},
\api{takewhile} to \code{TakeWhile}, \api{chain} to \code{Concat}, a
\api{groupby} to \code{GroupBy} po \emph{posortowanym} wejściu, i to ta
ostatnia zaskakuje ludzi.

\subsection{Sklejanie napisu na końcu}
\label{sec:ch05-join}

Nawyk z C\# na końcu potoku to \code{+=} w pętli, a folklor mówi, że w
Pythonie jest to kwadratowe. Na przypiętym interpreterze folklor myli się co
do przypadku, który ludzie faktycznie piszą, i ma rację co do przypadku,
którego nikt się nie spodziewa. Zapytaj interpretera: zdeasembluj rozgrzaną
pętlę, a kodem operacji jest \code{\valtext{ch05.concat.opcode}}. Gdy napis
po lewej stronie nie ma innej referencji, jest powiększany w miejscu, a
pętla jest liniowa.

\code{code/measure/ch05\_concat.py} podwaja pracę i sprawdza kształt, a nie
czas trwania, bo czas trwania jest własnością maszyny. Podwojenie wejścia
mnoży pracę liniową przez dwa, a kwadratową przez cztery, więc oba
ograniczenia leżą po dwóch stronach tej przerwy: zwykła pętla oraz
\code{"".join(...)} mieszczą się poniżej \val{ch05.concat.linear} przy
przejściu z \val{ch05.concat.small} doklejeń do \val{ch05.concat.large}, a
ta sama pętla z jedną dodatkową żywą referencją do napisu wypada powyżej
\val{ch05.concat.quadratic} przy przejściu z \val{ch05.concat.qsmall} do
\val{ch05.concat.qlarge}. Obie pary rozmiarów muszą się różnić i to jest ten
wynik: pętla kwadratowa przy rozmiarze liniowej biegnie minutami.

\begin{trapbox}
\textbf{Jest szybkie, dopóki niezwiązana linia nie uczyni go wolnym.}
Dopisanie bieżącego napisu do listy, domknięcie nad nim albo oddanie go
czemukolwiek, co go trzyma, unieważnia powiększanie w miejscu, a pętla robi
się kwadratowa bez żadnej zmiany w miejscu wywołania.
\code{"".join(pieces)} nie zależy od liczby referencji i nie ma złego
przypadku, i dlatego jest idiomem, chociaż \code{+=} zwykle wystarcza.
\end{trapbox}

\begin{exercise}{e05_03_pipeline}{To samo zapytanie, zachłanne i odroczone}
Napisz wersję posortowaną i zachłanną oraz odroczoną. Lista na poziomie
modułu zapisuje każdy element, którego potok dotknie, a testy z niej czytają,
by udowodnić, która z dwóch wersji już się wykonała, kiedy wracała \dash{}
więc buduj obie ze źródła, które dostajesz, a nie z listy zmaterializowanej
wcześniej. Ostatni test przechodzi po odroczonej wersji dwa razy.
\end{exercise}

\section{Generator to blok iteratora, i coś ponadto}
\label{sec:ch05-generators}

\code{yield} to \code{yield return}, a odwzorowanie sięga dalej, niż
większość ludzi się spodziewa: generator jest leniwy, wznawialny, a blok
\code{finally} w jego środku wykonuje się, gdy konsument przerwie wcześnie,
czyli robi to, co \code{IEnumerator.Dispose} na końcu \code{foreach}.

\pyregion{code/ch05/generators.py}{yield_return}{Blok iteratora, ze sprzątaniem przy wcześniejszym wyjściu włącznie}{lst:ch05-batched}

\noindent
A potem idzie dalej, w dwie rzeczy, których C\# nie ma jak zapisać.

\pyregion{code/ch05/generators.py}{two_way}{Wartość wysłana do środka i wartość zwrócona na zewnątrz}{lst:ch05-twoway}

\transcript{ch05-generators}

\mermaidfig{ch05-generator-resume}{Trzy sposoby wznowienia zawieszonego generatora i jedna rzecz, którą mają wspólną}{wznawianie generatora}

\begin{trapbox}
\textbf{\code{yield} jest wyrażeniem, nie instrukcją.} Ma wartość tego, co
konsument wysłał do środka \dash{} \code{None} przy zwykłym \code{next()} i
argument przy \code{send()}. Generator, który wykonuje \code{return}, wynosi
wartość także na zewnątrz, i dlatego \code{Generator} ma trzy parametry
typowe tam, gdzie \code{IEnumerable} ma jeden. To czyni pythonowy generator
korutyną w pierwotnym znaczeniu i jest to maszyneria, na której zbudowano
\code{async def}; rozdział~\ref{ch:asyncio} podejmuje rzecz w tym miejscu.
\end{trapbox}

\section{\code{with} to \code{using}, a contextlib to reszta}
\label{sec:ch05-with}

\code{using (var x = ...)} potrzebuje typu implementującego
\code{IDisposable}, czyli napisania klasy. \code{with} potrzebuje typu z
\code{\_\_enter\_\_} i \code{\_\_exit\_\_} \dash{} a
\api{contextlib.contextmanager} buduje taki typ z generatora, więc
typowym przypadkiem jest funkcja z \code{yield} pośrodku.

\pyregion{code/ch05/context.py}{contextmanager}{Wszystko przed yield to wejście; blok finally to wyjście}{lst:ch05-cm}

\begin{warning}
\code{try}/\code{finally} nie jest ozdobą. Bez niego wyjątek podniesiony w
ciele przechodzi przez \code{yield} na zewnątrz, a sprzątanie nigdy się nie
wykonuje \dash{} i to czyni tę rzecz jedyną tutaj, którą łatwiej zepsuć niż
instrukcję \code{using}, bo \code{using} nie da się napisać bez zwolnienia
zasobu.
\end{warning}

\noindent
Dwie kolejne różnice zarabiają na swoje miejsce i obie są w
\pkg{contextlib}:

\begin{itemize}[leftmargin=1.4em, itemsep=2pt]
  \item \api{suppress} \dash{} \code{\_\_exit\_\_} dostaje wyjątek i może
        go połknąć, czego \code{Dispose} nie potrafi.
  \item \api{ExitStack} \dash{} trzyma liczbę kontekstów nieznaną aż do
        czasu wykonania i rozwija je w odwrotnej kolejności, czego
        zagnieżdżone instrukcje \code{using} nie wyrażą.
\end{itemize}

\pyregion{code/ch05/context.py}{exitstack}{Liczba kontekstów rozstrzygnięta w czasie wykonania}{lst:ch05-exitstack}

\begin{versionbox}
Zapisz typ zwracany jako \code{Generator[str]}, a nie
\code{Iterator[str]}, który pokazuje każdy tutorial. Wtyczki typów, które
wysyła \pkg{pyright} \pyrightver{}, oznaczają to przeciążenie jako
przestarzałe, dosłownie: \enquote{annotating the return type as
\code{-> Iterator[Foo]} with \code{@contextmanager} is deprecated. Use
\code{-> Generator[Foo]} instead.} Kod działa tak i tak; zmieniła się
sprawdzarka.
\end{versionbox}

\begin{exercise}{e05_04_lease}{Menedżer kontekstu, który zawsze zwalnia}
Napisz \code{lease} przy użyciu \api{contextlib.contextmanager}. Testem,
który ma znaczenie, jest ten podnoszący wyjątek w ciele: menedżer kontekstu
oparty na generatorze bez \code{try}/\code{finally} wygląda poprawnie,
dopóki coś nie poleci. Czwarty test
zagnieżdża dwa i sprawdza, czy rozwijają się w odwrotnej kolejności.
\end{exercise}

\section{\code{match} i wzorzec, który nie jest porównaniem}
\label{sec:ch05-match}

C\#~9 i nowsze dają wyrażeniu switch wzorce relacyjne, typowe i
własnościowe. \code{match} ma je wszystkie, a do tego wzorce sekwencji i
odwzorowań oraz jedną regułę bez odpowiednika w C\#.

\pyregion{code/ch05/matching.py}{match}{Wzorce typu, własności, strażnika, sekwencji, odwzorowania i wartości}{lst:ch05-match}

\transcript{ch05-matching}

\noindent
Zwróć uwagę na zadeklarowaną unię nad funkcją. Parametr typu \code{object}
sprawia, że każde przechwycenie jest nieotypowane; unia jest tym, co ma
switch po zapieczętowanej hierarchii, i tym, co pozwala sprawdzarce zawężać
typ w każdym przypadku.

\begin{trapbox}
\textbf{Goła nazwa we wzorcu przechwytuje, a nie porównuje.}
\code{case ACTIVE:} nie sprawdza wartości względem stałej \code{ACTIVE}.
Dopasowuje się do czegokolwiek i przypisuje nazwę na nowo, lokalnie. Zapisem,
który porównuje, jest ten z kropką, \code{case Status.ACTIVE:}, a do reguły
dołączona jest jedna dobra wiadomość: Python odmawia kompilacji gołego
przechwycenia, po którym następuje kolejny przypadek, mówiąc, że pozostałe
wzorce są nieosiągalne. Kompiluje się po cichu tylko wtedy, gdy jest
przypadkiem ostatnim \dash{} czyli dokładnie tam, gdzie przeżywa przegląd
kodu.
\end{trapbox}

\pyregion{code/ch05/matching.py}{capture}{Jedyne miejsce, w którym reguła przechwycenia kompiluje się cicho}{lst:ch05-capture}

\begin{exercise}{e05_05_route}{Skieruj komunikat przy pomocy match}
Siedem przypadków, po kolei, obejmujących strażnika, wzorzec klasy ignorujący
pola, których nie nazywa, sekwencję, odwzorowanie i wariant zapasowy. Jeden
test podaje wartość, która pasowałaby do wzorca z gołą nazwą, i oczekuje, że
nie pasuje.
\end{exercise}

\section{Co to dało}
\label{sec:ch05-summary}

\begin{center}
\small
\begin{tabularx}{\linewidth}{@{}lX@{}}
\toprule
\textbf{C\#} & \textbf{Python} \\
\midrule
\code{Func<>}, delegat & funkcja z adnotacją \code{Callable[[A], B]} \\
atrybut czytany przez framework & nic; dekorator \emph{się wykonuje} \\
middleware wokół wywołania & dekorator \\
\code{params}, opcjonalne, nazwane & \code{*args}, domyślne, po nazwie \\
--- & znaczniki \code{/} i \code{*}, \code{**kwargs} \\
LINQ, odroczone i powtarzalne & składanie listy, zachłanne; albo wyrażenie
  generatora, odroczone i jednoprzebiegowe \\
\code{yield return} & \code{yield}, plus \code{send} i zwracana wartość \\
\code{using} / \code{IDisposable} & \code{with} / \api{contextmanager} \\
wyrażenie switch & \code{match}, ze wzorcami sekwencji i odwzorowań \\
\bottomrule
\end{tabularx}
\end{center}

\noindent
Zysk jest ten, od którego rozdział się zaczął. Wszystko powyżej obraca się
wokół tego, kiedy coś się wykonuje: przy definicji czy przy wywołaniu, przy
nawiasie czy przy przejściu, raz czy przy każdym przebiegu. Inżynier C\#
czytający Pythona szybko założy \emph{później} i pomyli się co do składania
list, a założy \emph{nigdy} i pomyli się co do dekoratorów. Pytaj o każdą
konstrukcję z tego rozdziału: czy to już się wydarzyło? Odpowiedź jest całą
różnicą.
